% Options for packages loaded elsewhere
% Options for packages loaded elsewhere
\PassOptionsToPackage{unicode}{hyperref}
\PassOptionsToPackage{hyphens}{url}
\PassOptionsToPackage{dvipsnames,svgnames,x11names}{xcolor}
%
\documentclass[
  11pt,
  letterpaper,
  DIV=11,
  numbers=noendperiod]{scrartcl}
\usepackage{xcolor}
\usepackage[margin=2.5cm]{geometry}
\usepackage{amsmath,amssymb}
\setcounter{secnumdepth}{5}
\usepackage{iftex}
\ifPDFTeX
  \usepackage[T1]{fontenc}
  \usepackage[utf8]{inputenc}
  \usepackage{textcomp} % provide euro and other symbols
\else % if luatex or xetex
  \usepackage{unicode-math} % this also loads fontspec
  \defaultfontfeatures{Scale=MatchLowercase}
  \defaultfontfeatures[\rmfamily]{Ligatures=TeX,Scale=1}
\fi
\usepackage{lmodern}
\ifPDFTeX\else
  % xetex/luatex font selection
\fi
% Use upquote if available, for straight quotes in verbatim environments
\IfFileExists{upquote.sty}{\usepackage{upquote}}{}
\IfFileExists{microtype.sty}{% use microtype if available
  \usepackage[]{microtype}
  \UseMicrotypeSet[protrusion]{basicmath} % disable protrusion for tt fonts
}{}
\makeatletter
\@ifundefined{KOMAClassName}{% if non-KOMA class
  \IfFileExists{parskip.sty}{%
    \usepackage{parskip}
  }{% else
    \setlength{\parindent}{0pt}
    \setlength{\parskip}{6pt plus 2pt minus 1pt}}
}{% if KOMA class
  \KOMAoptions{parskip=half}}
\makeatother
% Make \paragraph and \subparagraph free-standing
\makeatletter
\ifx\paragraph\undefined\else
  \let\oldparagraph\paragraph
  \renewcommand{\paragraph}{
    \@ifstar
      \xxxParagraphStar
      \xxxParagraphNoStar
  }
  \newcommand{\xxxParagraphStar}[1]{\oldparagraph*{#1}\mbox{}}
  \newcommand{\xxxParagraphNoStar}[1]{\oldparagraph{#1}\mbox{}}
\fi
\ifx\subparagraph\undefined\else
  \let\oldsubparagraph\subparagraph
  \renewcommand{\subparagraph}{
    \@ifstar
      \xxxSubParagraphStar
      \xxxSubParagraphNoStar
  }
  \newcommand{\xxxSubParagraphStar}[1]{\oldsubparagraph*{#1}\mbox{}}
  \newcommand{\xxxSubParagraphNoStar}[1]{\oldsubparagraph{#1}\mbox{}}
\fi
\makeatother

\usepackage{color}
\usepackage{fancyvrb}
\newcommand{\VerbBar}{|}
\newcommand{\VERB}{\Verb[commandchars=\\\{\}]}
\DefineVerbatimEnvironment{Highlighting}{Verbatim}{commandchars=\\\{\}}
% Add ',fontsize=\small' for more characters per line
\usepackage{framed}
\definecolor{shadecolor}{RGB}{241,243,245}
\newenvironment{Shaded}{\begin{snugshade}}{\end{snugshade}}
\newcommand{\AlertTok}[1]{\textcolor[rgb]{0.68,0.00,0.00}{#1}}
\newcommand{\AnnotationTok}[1]{\textcolor[rgb]{0.37,0.37,0.37}{#1}}
\newcommand{\AttributeTok}[1]{\textcolor[rgb]{0.40,0.45,0.13}{#1}}
\newcommand{\BaseNTok}[1]{\textcolor[rgb]{0.68,0.00,0.00}{#1}}
\newcommand{\BuiltInTok}[1]{\textcolor[rgb]{0.00,0.23,0.31}{#1}}
\newcommand{\CharTok}[1]{\textcolor[rgb]{0.13,0.47,0.30}{#1}}
\newcommand{\CommentTok}[1]{\textcolor[rgb]{0.37,0.37,0.37}{#1}}
\newcommand{\CommentVarTok}[1]{\textcolor[rgb]{0.37,0.37,0.37}{\textit{#1}}}
\newcommand{\ConstantTok}[1]{\textcolor[rgb]{0.56,0.35,0.01}{#1}}
\newcommand{\ControlFlowTok}[1]{\textcolor[rgb]{0.00,0.23,0.31}{\textbf{#1}}}
\newcommand{\DataTypeTok}[1]{\textcolor[rgb]{0.68,0.00,0.00}{#1}}
\newcommand{\DecValTok}[1]{\textcolor[rgb]{0.68,0.00,0.00}{#1}}
\newcommand{\DocumentationTok}[1]{\textcolor[rgb]{0.37,0.37,0.37}{\textit{#1}}}
\newcommand{\ErrorTok}[1]{\textcolor[rgb]{0.68,0.00,0.00}{#1}}
\newcommand{\ExtensionTok}[1]{\textcolor[rgb]{0.00,0.23,0.31}{#1}}
\newcommand{\FloatTok}[1]{\textcolor[rgb]{0.68,0.00,0.00}{#1}}
\newcommand{\FunctionTok}[1]{\textcolor[rgb]{0.28,0.35,0.67}{#1}}
\newcommand{\ImportTok}[1]{\textcolor[rgb]{0.00,0.46,0.62}{#1}}
\newcommand{\InformationTok}[1]{\textcolor[rgb]{0.37,0.37,0.37}{#1}}
\newcommand{\KeywordTok}[1]{\textcolor[rgb]{0.00,0.23,0.31}{\textbf{#1}}}
\newcommand{\NormalTok}[1]{\textcolor[rgb]{0.00,0.23,0.31}{#1}}
\newcommand{\OperatorTok}[1]{\textcolor[rgb]{0.37,0.37,0.37}{#1}}
\newcommand{\OtherTok}[1]{\textcolor[rgb]{0.00,0.23,0.31}{#1}}
\newcommand{\PreprocessorTok}[1]{\textcolor[rgb]{0.68,0.00,0.00}{#1}}
\newcommand{\RegionMarkerTok}[1]{\textcolor[rgb]{0.00,0.23,0.31}{#1}}
\newcommand{\SpecialCharTok}[1]{\textcolor[rgb]{0.37,0.37,0.37}{#1}}
\newcommand{\SpecialStringTok}[1]{\textcolor[rgb]{0.13,0.47,0.30}{#1}}
\newcommand{\StringTok}[1]{\textcolor[rgb]{0.13,0.47,0.30}{#1}}
\newcommand{\VariableTok}[1]{\textcolor[rgb]{0.07,0.07,0.07}{#1}}
\newcommand{\VerbatimStringTok}[1]{\textcolor[rgb]{0.13,0.47,0.30}{#1}}
\newcommand{\WarningTok}[1]{\textcolor[rgb]{0.37,0.37,0.37}{\textit{#1}}}

\usepackage{longtable,booktabs,array}
\usepackage{calc} % for calculating minipage widths
% Correct order of tables after \paragraph or \subparagraph
\usepackage{etoolbox}
\makeatletter
\patchcmd\longtable{\par}{\if@noskipsec\mbox{}\fi\par}{}{}
\makeatother
% Allow footnotes in longtable head/foot
\IfFileExists{footnotehyper.sty}{\usepackage{footnotehyper}}{\usepackage{footnote}}
\makesavenoteenv{longtable}
\usepackage{graphicx}
\makeatletter
\newsavebox\pandoc@box
\newcommand*\pandocbounded[1]{% scales image to fit in text height/width
  \sbox\pandoc@box{#1}%
  \Gscale@div\@tempa{\textheight}{\dimexpr\ht\pandoc@box+\dp\pandoc@box\relax}%
  \Gscale@div\@tempb{\linewidth}{\wd\pandoc@box}%
  \ifdim\@tempb\p@<\@tempa\p@\let\@tempa\@tempb\fi% select the smaller of both
  \ifdim\@tempa\p@<\p@\scalebox{\@tempa}{\usebox\pandoc@box}%
  \else\usebox{\pandoc@box}%
  \fi%
}
% Set default figure placement to htbp
\def\fps@figure{htbp}
\makeatother





\setlength{\emergencystretch}{3em} % prevent overfull lines

\providecommand{\tightlist}{%
  \setlength{\itemsep}{0pt}\setlength{\parskip}{0pt}}



 


\KOMAoption{captions}{tableheading}
\makeatletter
\@ifpackageloaded{caption}{}{\usepackage{caption}}
\AtBeginDocument{%
\ifdefined\contentsname
  \renewcommand*\contentsname{Table of contents}
\else
  \newcommand\contentsname{Table of contents}
\fi
\ifdefined\listfigurename
  \renewcommand*\listfigurename{List of Figures}
\else
  \newcommand\listfigurename{List of Figures}
\fi
\ifdefined\listtablename
  \renewcommand*\listtablename{List of Tables}
\else
  \newcommand\listtablename{List of Tables}
\fi
\ifdefined\figurename
  \renewcommand*\figurename{Figure}
\else
  \newcommand\figurename{Figure}
\fi
\ifdefined\tablename
  \renewcommand*\tablename{Table}
\else
  \newcommand\tablename{Table}
\fi
}
\@ifpackageloaded{float}{}{\usepackage{float}}
\floatstyle{ruled}
\@ifundefined{c@chapter}{\newfloat{codelisting}{h}{lop}}{\newfloat{codelisting}{h}{lop}[chapter]}
\floatname{codelisting}{Listing}
\newcommand*\listoflistings{\listof{codelisting}{List of Listings}}
\makeatother
\makeatletter
\makeatother
\makeatletter
\@ifpackageloaded{caption}{}{\usepackage{caption}}
\@ifpackageloaded{subcaption}{}{\usepackage{subcaption}}
\makeatother
\usepackage{bookmark}
\IfFileExists{xurl.sty}{\usepackage{xurl}}{} % add URL line breaks if available
\urlstyle{same}
\hypersetup{
  pdftitle={Fengping River Hydropower --- Test},
  pdfauthor={{[}MAYA SCANLON , WILLIAM LIEN{]}},
  colorlinks=true,
  linkcolor={blue},
  filecolor={Maroon},
  citecolor={Blue},
  urlcolor={Blue},
  pdfcreator={LaTeX via pandoc}}


\title{Fengping River Hydropower --- Test}
\usepackage{etoolbox}
\makeatletter
\providecommand{\subtitle}[1]{% add subtitle to \maketitle
  \apptocmd{\@title}{\par {\large #1 \par}}{}{}
}
\makeatother
\subtitle{Assignment 6 --- 3 simple function tests}
\author{{[}MAYA SCANLON , WILLIAM LIEN{]}}
\date{2026-06-14}
\begin{document}
\maketitle

\renewcommand*\contentsname{Table of contents}
{
\hypersetup{linkcolor=}
\setcounter{tocdepth}{2}
\tableofcontents
}

This document is a newly organised document on the test cases for the
three functions we have created in the module. The tests checks the
input and output format. The tests are written using the
\texttt{testthat} package, which provides a framework for writing unit
tests in R. Each function is tested for its expected behavior, including
handling of edge cases and invalid inputs.

\begin{Shaded}
\begin{Highlighting}[]
\FunctionTok{library}\NormalTok{(tidyverse)}
\FunctionTok{library}\NormalTok{(here)}
\FunctionTok{library}\NormalTok{(testthat)}

\FunctionTok{source}\NormalTok{(}\FunctionTok{here}\NormalTok{(}\StringTok{"module"}\NormalTok{, }\StringTok{"hydro\_reservoir.R"}\NormalTok{))}
\FunctionTok{source}\NormalTok{(}\FunctionTok{here}\NormalTok{(}\StringTok{"module"}\NormalTok{, }\StringTok{"hydro\_sediment.R"}\NormalTok{))}
\FunctionTok{source}\NormalTok{(}\FunctionTok{here}\NormalTok{(}\StringTok{"module"}\NormalTok{, }\StringTok{"hydro\_finance.R"}\NormalTok{))}
\end{Highlighting}
\end{Shaded}

\section{\texorpdfstring{Function 1: \texttt{fill\_daily\_na()} ---
daily flow
data}{Function 1: fill\_daily\_na() --- daily flow data}}\label{function-1-fill_daily_na-daily-flow-data}

\begin{Shaded}
\begin{Highlighting}[]
\FunctionTok{test\_that}\NormalTok{(}\StringTok{"fill\_daily\_na: detects and fills NA gaps, output is reasonable"}\NormalTok{, \{}

  \CommentTok{\# Synthetic flow data with a gap (NA) in the middle}
\NormalTok{  synth }\OtherTok{\textless{}{-}} \FunctionTok{data.frame}\NormalTok{(}
    \AttributeTok{date  =} \FunctionTok{seq}\NormalTok{(}\FunctionTok{as.Date}\NormalTok{(}\StringTok{"2020{-}01{-}01"}\NormalTok{), }\FunctionTok{as.Date}\NormalTok{(}\StringTok{"2020{-}01{-}05"}\NormalTok{), }\AttributeTok{by =} \StringTok{"day"}\NormalTok{),}
    \AttributeTok{Q\_cms =} \FunctionTok{c}\NormalTok{(}\DecValTok{10}\NormalTok{, }\DecValTok{12}\NormalTok{, }\ConstantTok{NA}\NormalTok{, }\DecValTok{16}\NormalTok{, }\DecValTok{18}\NormalTok{)}
\NormalTok{  )}

  \CommentTok{\# Confirm the test data actually has a gap}
  \FunctionTok{expect\_true}\NormalTok{(}\FunctionTok{any}\NormalTok{(}\FunctionTok{is.na}\NormalTok{(synth}\SpecialCharTok{$}\NormalTok{Q\_cms)))}

\NormalTok{  result }\OtherTok{\textless{}{-}} \FunctionTok{fill\_daily\_na}\NormalTok{(synth, }\AttributeTok{method =} \StringTok{"linear"}\NormalTok{)}

  \CommentTok{\# Gap is gone after filling}
  \FunctionTok{expect\_true}\NormalTok{(}\FunctionTok{all}\NormalTok{(}\SpecialCharTok{!}\FunctionTok{is.na}\NormalTok{(result}\SpecialCharTok{$}\NormalTok{Q\_cms)))}

  \CommentTok{\# Filled flow values are physically reasonable (non{-}negative)}
  \FunctionTok{expect\_true}\NormalTok{(}\FunctionTok{all}\NormalTok{(result}\SpecialCharTok{$}\NormalTok{Q\_cms }\SpecialCharTok{\textgreater{}=} \DecValTok{0}\NormalTok{))}

  \CommentTok{\# Overall distribution of filled values}
  \FunctionTok{summary}\NormalTok{(result}\SpecialCharTok{$}\NormalTok{Q\_cms)}
\NormalTok{\})}
\end{Highlighting}
\end{Shaded}

\begin{verbatim}
Test passed with 3 successes.
\end{verbatim}

\section{\texorpdfstring{Function 2: \texttt{compute\_npv()} --- NPV /
CAPEX}{Function 2: compute\_npv() --- NPV / CAPEX}}\label{function-2-compute_npv-npv-capex}

\begin{Shaded}
\begin{Highlighting}[]
\FunctionTok{test\_that}\NormalTok{(}\StringTok{"compute\_npv: basic correctness and input checks"}\NormalTok{, \{}

  \CommentTok{\# At r = 0, NPV is just the sum of cash flows (hand{-}checkable)}
\NormalTok{  cf }\OtherTok{\textless{}{-}} \FunctionTok{c}\NormalTok{(}\SpecialCharTok{{-}}\DecValTok{1000}\NormalTok{, }\DecValTok{300}\NormalTok{, }\DecValTok{300}\NormalTok{, }\DecValTok{300}\NormalTok{, }\DecValTok{300}\NormalTok{)}
  \FunctionTok{expect\_equal}\NormalTok{(}\FunctionTok{compute\_npv}\NormalTok{(cf, }\AttributeTok{discount\_rate =} \DecValTok{0}\NormalTok{), }\FunctionTok{sum}\NormalTok{(cf))}

  \CommentTok{\# A single positive value at year 0 needs no discounting}
  \FunctionTok{expect\_equal}\NormalTok{(}\FunctionTok{compute\_npv}\NormalTok{(}\FunctionTok{c}\NormalTok{(}\DecValTok{100}\NormalTok{, }\DecValTok{0}\NormalTok{), }\AttributeTok{discount\_rate =} \FloatTok{0.1}\NormalTok{), }\DecValTok{100}\NormalTok{)}

  \CommentTok{\# Higher discount rate {-}\textgreater{} lower NPV (for positive future flows)}
\NormalTok{  npv\_low  }\OtherTok{\textless{}{-}} \FunctionTok{compute\_npv}\NormalTok{(cf, }\AttributeTok{discount\_rate =} \FloatTok{0.01}\NormalTok{)}
\NormalTok{  npv\_high }\OtherTok{\textless{}{-}} \FunctionTok{compute\_npv}\NormalTok{(cf, }\AttributeTok{discount\_rate =} \FloatTok{0.15}\NormalTok{)}
  \FunctionTok{expect\_true}\NormalTok{(npv\_low }\SpecialCharTok{\textgreater{}}\NormalTok{ npv\_high)}

  \CommentTok{\# Invalid input (length \textless{} 2) should error, per stopifnot()}
  \FunctionTok{expect\_error}\NormalTok{(}\FunctionTok{compute\_npv}\NormalTok{(}\FunctionTok{c}\NormalTok{(}\SpecialCharTok{{-}}\DecValTok{1000}\NormalTok{), }\AttributeTok{discount\_rate =} \FloatTok{0.05}\NormalTok{))}
\NormalTok{\})}
\end{Highlighting}
\end{Shaded}

\begin{verbatim}
Test passed with 4 successes.
\end{verbatim}

\section{\texorpdfstring{Function 3: \texttt{fit\_rating\_curve()} ---
sediment data format \&
fit}{Function 3: fit\_rating\_curve() --- sediment data format \& fit}}\label{function-3-fit_rating_curve-sediment-data-format-fit}

\begin{Shaded}
\begin{Highlighting}[]
\FunctionTok{test\_that}\NormalTok{(}\StringTok{"fit\_rating\_curve: input format and fitted curve are valid"}\NormalTok{, \{}

  \CommentTok{\# Sediment sampling data: required columns Discharge\_CMS, Sus\_Load\_MTDay, flag}
\NormalTok{  sed\_sample }\OtherTok{\textless{}{-}} \FunctionTok{data.frame}\NormalTok{(}
    \AttributeTok{Discharge\_CMS  =} \FunctionTok{c}\NormalTok{(}\DecValTok{10}\NormalTok{, }\DecValTok{20}\NormalTok{, }\DecValTok{30}\NormalTok{, }\DecValTok{50}\NormalTok{, }\DecValTok{80}\NormalTok{, }\DecValTok{120}\NormalTok{),}
    \AttributeTok{Sus\_Load\_MTDay =} \FunctionTok{c}\NormalTok{(}\DecValTok{50}\NormalTok{, }\DecValTok{150}\NormalTok{, }\DecValTok{300}\NormalTok{, }\DecValTok{700}\NormalTok{, }\DecValTok{1500}\NormalTok{, }\DecValTok{3000}\NormalTok{),}
    \AttributeTok{flag           =} \StringTok{"ok"}
\NormalTok{  )}

  \CommentTok{\# Required columns are present}
  \FunctionTok{expect\_true}\NormalTok{(}\FunctionTok{all}\NormalTok{(}\FunctionTok{c}\NormalTok{(}\StringTok{"Discharge\_CMS"}\NormalTok{, }\StringTok{"Sus\_Load\_MTDay"}\NormalTok{, }\StringTok{"flag"}\NormalTok{) }\SpecialCharTok{\%in\%} \FunctionTok{names}\NormalTok{(sed\_sample)))}

  \CommentTok{\# At least 5 valid rows for the regression}
\NormalTok{  valid\_rows }\OtherTok{\textless{}{-}}\NormalTok{ sed\_sample}\SpecialCharTok{$}\NormalTok{Discharge\_CMS }\SpecialCharTok{\textgreater{}} \DecValTok{0} \SpecialCharTok{\&}\NormalTok{ sed\_sample}\SpecialCharTok{$}\NormalTok{Sus\_Load\_MTDay }\SpecialCharTok{\textgreater{}} \DecValTok{0}
  \FunctionTok{expect\_true}\NormalTok{(}\FunctionTok{sum}\NormalTok{(valid\_rows) }\SpecialCharTok{\textgreater{}=} \DecValTok{5}\NormalTok{)}

  \CommentTok{\# Fitted curve returns a list with \textquotesingle{}a\textquotesingle{} and \textquotesingle{}b\textquotesingle{}}
\NormalTok{  rc }\OtherTok{\textless{}{-}} \FunctionTok{fit\_rating\_curve}\NormalTok{(sed\_sample, }\AttributeTok{plot\_fit =} \ConstantTok{FALSE}\NormalTok{)}
  \FunctionTok{expect\_true}\NormalTok{(}\FunctionTok{all}\NormalTok{(}\FunctionTok{c}\NormalTok{(}\StringTok{"a"}\NormalTok{, }\StringTok{"b"}\NormalTok{) }\SpecialCharTok{\%in\%} \FunctionTok{names}\NormalTok{(rc)))}

  \CommentTok{\# \textquotesingle{}b\textquotesingle{} (power{-}law exponent) is positive: sediment load increases with flow}
  \FunctionTok{expect\_true}\NormalTok{(rc}\SpecialCharTok{$}\NormalTok{b }\SpecialCharTok{\textgreater{}} \DecValTok{0}\NormalTok{)}
\NormalTok{\})}
\end{Highlighting}
\end{Shaded}

\begin{verbatim}
Test passed with 4 successes.
\end{verbatim}




\end{document}
